\documentclass[10pt,a4paper]{amsart}
\usepackage[margin=19mm]{geometry}
\usepackage{amsmath,amssymb,mathtools}
\usepackage[scr=rsfs,cal=euler]{mathalfa}
\usepackage{xfrac}
\usepackage[unicode,hidelinks,bookmarks=false]{hyperref}
\newcommand{\ie}{i.e.\ }
\newcommand{\cf}{cf.\ }
\newcommand{\resp}{respectively\ }
\newcommand{\Subsection}{\subsection}
\providecommand{\nobreakdash}{\nobreak\hbox{-}}
\setlength{\emergencystretch}{2em}
\multlinegap0pt
\usepackage{amssymb}
\usepackage{tcolorbox}
\usepackage{dsfont}
\usepackage{mathtools}
\usepackage{float}
\usepackage{enumerate}
\usepackage{tikz}
\usepackage[shortlabels]{enumitem}
\usepackage{stackengine}
\usepackage{xcolor}
\usepackage[normalem]{ulem}
\newtheorem{definition}{Definition}
\newtheorem{remark}{Remark}
\newtheorem{theorem}{Theorem}
\newtheorem{lemma}{Lemma}
\usepackage{cleveref}
\crefname{definition}{Definition}{Definitions}
\crefname{remark}{Remark}{Remarks}
\crefname{theorem}{Theorem}{Theorems}
\crefname{lemma}{Lemma}{Lemmas}
\newcommand{\B}{\textcolor{orange}}
\def\T{\mathcal{T}}
\newcommand{\di}[1]{\,\mathrm{d}#1}
\def\dzdy{\di{z}\di{y}\di{\tau}\di{\overline{x}}\di{t}}
\def\es{\varepsilon}
\def\mint{\int_{0}^{T}\int_{\Sigma}\int_{\mathcal{T}}\int_{Y}\int_{Z}}
\def\oii{\Omega_i^\varepsilon}
\def\oll{\Omega_L^\varepsilon}
\def\omm{\Omega_M^\varepsilon}
\def\orr{\Omega_R^\varepsilon}
\def\p{\partial}
\title{Interface Conditions for Wave Propagation Through a Time-varying Metasurface}
\author{Vishnu Raveendran$^a$\, and Barbara Verfürth$^a$}
\date{Source: arXiv:2608.27061v1 (CC BY 4.0)}
\begin{document}
\maketitle

\noindent\emph{Source and licence.} Interface Conditions for Wave Propagation Through a Time-varying Metasurface. Version arXiv:2608.27061v1 (CC BY 4.0), distributed by arXiv under the Creative Commons Attribution 4.0 International licence, http://creativecommons.org/licenses/by/4.0/, as recorded in the arXiv licence metadata for this exact version and shown on the abstract page https://arxiv.org/abs/2608.27061v1. Full licence notice: \texttt{licenses/arxiv-2608.27061-CC-BY-4.0.txt}.\par\smallskip

\noindent\textit{Editorial scope.} Counted unit: the article's lemma eq:strong\_form\_u\_1 (source0.tex lines 991--1013). The article's complete original proof of it is delivered with it (source0.tex lines 1014--1214, one \texttt{proof} environment of 201 lines). No mathematics is written, repaired or re-derived: the statement and the proof are the article's own lines, in the article's own notation, and no retained line is substituted or re-flowed.  Retained uncounted as the source-local context the proof needs: lines 377--382; lines 419--423; lines 701--713; lines 813--823; lines 816--820.  Nothing was omitted from any retained span, so the package carries no omission record.  No retained line contains a citation, so the wrapper carries no bibliography and no bibliography entry.  No retained line contains a figure environment, an image-inclusion command or drawing code, so no image is delivered and the package declares no asset dependency.  Editorial formatting only, in addition: an AMS \texttt{amsart} preamble that restates the article's own declarations and macros used by the retained lines, namely the environments \texttt{definition}, \texttt{remark}, \texttt{theorem}, \texttt{lemma} on separate counters, together with the standard packages those lines need.  The retained environments print in this document as Definition 1, Remark 1, Theorem 1, Lemma 1 in order of appearance, whereas the article prints the same items with the numbering of its own full document.  The complete native source of the article as distributed in the arXiv e-print for this exact version is bundled unchanged as \texttt{sources/208681/source0.tex}.
\begin{multline}\label{eq:weakform}
    -\rho_L\int_0^T\int_{\oll}  \dfrac{\partial u_L^\es}{\partial t}  \dfrac{\partial \phi_L}{\partial t}\di{x}\di{t}-  \dfrac{1}{\es} \int_0^T\int_{\omm}  \left({\rho_M^\es} \dfrac{\partial u_M^\es}{\partial t}\right) \dfrac{\partial \phi_M}{\partial t}\di{x}\di{t} -    \rho_R\int_0^T\int_{\orr} \dfrac{\partial u_R^\es}{\partial t}\dfrac{\partial \phi_R}{\partial t} \di{x}\di{t}\\
    +D_L\int_0^T\int_{\oll}\nabla u_L^\es \cdot \nabla \phi_L   \di{x}\di{t} +   \dfrac{1}{\es}\int_0^T\int_{\omm}  D_M^\es  \nabla u_M^\es\cdot\nabla \phi_M \di{x}\di{t}+D_R\int_0^T\int_{\orr} \nabla u_R^\es\cdot\nabla \phi_R \di{x} \di{t}
   \\
    = \int_0^T\int_{\oll} f_L\phi_L \di{x}\di{t}+\int_{\oll}\rho_L v_2\phi_L(0)\di{x}
\end{multline}

	\begin{align}
     \left\|\ u_L^\es\right\|_{L^{\infty}(0,T;L^2(\oll))}+  \left\|\ u_R^\es\right\|_{L^{\infty}(0,T;L^2(\orr))}+  \dfrac{1}{\sqrt{\es}}\left\|\ u_M^\es\right\|_{L^{\infty}(0,T;L^2(\omm))}&\leq C,\label{eq:energy_estimate_1_case_1}\\
	    \left\|\ \dfrac{\partial u_L^\es}{\partial t}\right\|_{L^{\infty}(0,T;L^2(\oll))}+  \left\|\ \dfrac{\partial u_R^\es}{\partial t}\right\|_{L^{\infty}(0,T;L^2(\orr))}+  \dfrac{1}{\sqrt{\es}}\left\|\ \dfrac{\partial u_M^\es}{\partial t}\right\|_{L^{\infty}(0,T;L^2(\omm))}&\leq C,\label{eq:energy_estimate_2_case_1}\\
	    \|\nabla u_L^\es\|_{L^{\infty}(0,T;L^2(\oll))}+\|\nabla u_R^\es\|_{L^{\infty}(0,T;L^2(\orr))}+\dfrac{1}{\sqrt{\es}}\|\nabla u_M^\es\|_{L^{\infty}(0,T;L^2(\omm))}&\leq C\label{eq:energy_estimate_3_case_1},
	\end{align}

	\begin{definition}\label{def:multiscale_convergence}
A sequence $v^\es\in L^2(0,T;L^2(\omm))$ is said to be multi-scale convergent (with three spatial and two temporal scales) to $v^0(t,\overline{x},\tau,y,z)\in L^2((0,T)\times \Sigma\times  \T \times Y\times Z)$, if
\begin{multline*}
    \lim_{\es\rightarrow 0}\dfrac{1}{\es}\int_0^T\int_{\omm}v^\es(t,x)\phi\left(t,\overline{x},\dfrac{t}{\es},\dfrac{x}{\es},\dfrac{x_2}{\es^2}\right)\di{t}\di{x}\\
=\mint v^0(t,\overline{x},\tau,y,z)\phi(t,\overline{x},\tau,y,z)\dzdy
\end{multline*}
for all $\phi \in L^2((0,T)\times \Sigma;\mathcal{C}_{\sharp}(\overline{\T\times Y\times Z}))$.
\newline

\begin{remark}
    The notion of multiscale convergence for thin layers can be extended naturally to an arbitrary number of well-separated spatial and temporal scales.
\end{remark}
	\end{definition}

   \begin{theorem}\label{theorem:convergence_in_the_layer}
    Assume $(A1)-(A2)$ and \eqref{eq:energy_estimate_1_case_1}--\eqref{eq:energy_estimate_3_case_1} hold. Then there exist \\ $u_M^0\in L^2((0,T);H^1_{\sharp}(\Sigma))\cap H^1((0,T);L^2(\Sigma))$, $u_M^1\in L^2((0,T)\times \Sigma;\mathcal{H}_{\sharp,0}(\T\times Y))$ and $u_M^2\in L^2((0,T)\times \Sigma \times \T\times Y;H^1_{\sharp,0}(Z) )$ such that up to a subsequence
    \begin{subequations}
    \begin{align}
        u_M^\es(t,x)&\longrightarrow u_M^0(t,\overline{x})\label{eq:multiscale_convergence_for_u}\\
        \nabla u_M^\es(t,x)&\longrightarrow \nabla_{\overline{x}}u_M^0(t,\overline{x})+\nabla_y u_M^1(t,\overline{x},\tau,y)+\frac{d}{dz}u_M^2(t,\overline{x},\tau,y,z)e_2\label{eq:multiscale_convergence_for_gradient_term} \\
        \dfrac{\partial u_M^\es(t,x)}{\partial t}&\longrightarrow  \dfrac{\partial u_M^0(t,x)}{\partial t}+ \dfrac{\partial u_M^1(t,\overline{x},\tau,y)}{\partial \tau}\label{eq:multiscale_convergence_for_time_derivative_term}
    \end{align}
    \end{subequations}
    where the convergences are in the sense of multi-scale convergence for the thin layer.
    \end{theorem}

    \begin{align}
        u_M^\es(t,x)&\longrightarrow u_M^0(t,\overline{x})\\
        \nabla u_M^\es(t,x)&\longrightarrow \nabla_{\overline{x}}u_M^0(t,\overline{x})+\nabla_y u_M^1(t,\overline{x},\tau,y)+\frac{d}{dz}u_M^2(t,\overline{x},\tau,y,z)e_2 \\
        \dfrac{\partial u_M^\es(t,x)}{\partial t}&\longrightarrow  \dfrac{\partial u_M^0(t,x)}{\partial t}+ \dfrac{\partial u_M^1(t,\overline{x},\tau,y)}{\partial \tau}
    \end{align}

    \begin{lemma}
   Assume $(A1)-(A2)$ and \eqref{eq:energy_estimate_1_case_1}--\eqref{eq:energy_estimate_3_case_1} hold. Then, in the distributional sense, the functions $u_M^0,u_M^1$ and $u_M^2$ from  \Cref{theorem:convergence_in_the_layer} satisfy
 \begin{subequations}  

\begin{multline}\label{eq:strong_form_u_1}
   \dfrac{\p}{\p \tau} \int_Z \left( {\rho_M(t,\overline{x},\tau,y,z)} \left(\dfrac{\partial u_M^0(t,x)}{\partial t}+ \dfrac{\partial u_M^1(t,\overline{x},\tau,y)}{\partial \tau}\right)\right)\di{z}\\
     -\mathrm{div_y}\int_Z \bigg( D_M(t,\overline{x},\tau,y,z)
    \left( \nabla_{\overline{x}}u_M^0(t,\overline{x})+\nabla_y u_M^1(t,\overline{x},\tau,y)+\nabla_zu_M^2(t,\overline{x},\tau,y,z)\right)\bigg)\di{z}=0\\
    \mbox{for a.e.} (t,\overline{x},\tau,y)\in (0,T)\times\Sigma\times (0,1)\times Y,
\end{multline}
along with the boundary condition
\begin{align}\label{eq:strong_form_u_1_bc}
    \bigg( D_M(t,\overline{x},\tau,y,z)
    \left( \nabla_{\overline{x}}u_M^0(t,\overline{x})+\nabla_y u_M^1(t,\overline{x},\tau,y)+\nabla_zu_M^2(t,\overline{x},\tau,y,z)\right)\bigg)\cdot n=0\quad\mbox{on}\quad Y_L\cup Y_R,
\end{align}
and
\begin{multline}\label{eq:strong_form_for_u_2}
    -\mathrm{div_z}\Big(   D_M(t,\overline{x},\tau,y,z)
    \left( \nabla_{\overline{x}}u_M^0(t,\overline{x})+\nabla_y u_M^1(t,\overline{x},\tau,y)+\nabla_zu_M^2(t,\overline{x},\tau,y,z)\right)\Big)=0\\
    \mbox{for a.e. } (t,\overline{x},\tau,y,z)\in (0,T)\times\Sigma\times (0,1)\times Y\times Z.
\end{multline}
\end{subequations}
\end{lemma}

\begin{proof}
    We first construct a test function in $(0,T)\times\Omega$ with suitable regularity and support in both fast and slow variables, which is then used in the weak formulation of the microscopic problem. Then, using the uniform energy estimates and the multiscale convergence for thin layer, as $\es\rightarrow0$, we deduce \eqref{eq:strong_form_u_1} and \eqref{eq:strong_form_u_1_bc}.
Let
\begin{align*}
    \phi_1:=\begin{cases} 
      \alpha_1(t)\alpha_2(\frac{t}{\es})\beta(\bar x) \psi_L(\frac{\bar x}{\es})\rho(\frac{x_1}{\es}) &  \qquad \mbox{in}\qquad (0,T)\times \oll\\
      \alpha_1(t)\alpha_2(\frac{t}{\es})\beta(\bar x) \psi_M(\frac{x}{\es}) & \qquad \mbox{in}\qquad (0,T)\times \omm\\
      \alpha_1(t)\alpha_2(\frac{t}{\es})\beta(\bar x) \psi_R(\frac{\bar x}{\es})\rho(\frac{-x_1}{\es}) & \qquad \mbox{in}\qquad (0,T)\times \orr,
   \end{cases}
\end{align*}
where
\begin{itemize}
    \item $\alpha_1\in C_c^\infty((0,T))$ and $\beta\in C_c^\infty(\Sigma).$
    \item $ \alpha_2\in C_{c}^\infty([0,1])$ which extends periodically to whole $\mathbb{R}.$
    \item $\psi_M\in C_c^\infty(\overline Y) $ which extends periodically with respect to $y_2$ variable.
    \item For $i\in \{L,R\}$, $ \psi_i $ is defined in $(0,T)\times \overline{\oii}$ such that $\psi_i\equiv\psi_M$ on $i^\es$ and independent of $x_1$ variable.
    \item $\rho\in C^\infty([0,\infty))$ is a cut off function satisfying  $0\leq\rho \leq1$, $\rho=1$ in $[0,1]$ and $\rho=0 $ in $[2,\infty).$
\end{itemize}







% where $\alpha_1\in C_c^\infty((0,T))$,  $ \alpha_2\in C_{c}^\infty([0,1])$ and extent periodically to whole $\mathbb{R},$ $\beta\in C_c^\infty(\Sigma)$, $\psi_M\in C_c^\infty(\overline Y) $ and extend periodically with respect to $y_2$ variable, $\psi_L $ defined in $(0,T)\times \overline{\oll}$ such that $\psi_L\equiv\psi_M$ on $L^\es$ and independent on $x_1$ variable, similarly we define $\psi_R $ defined in $(0,T)\times \overline{\orr}$ such that $\psi_R\equiv\psi_M$ on $R^\es$ and independent on $x_1$ variable, $\rho\in C^\infty([0,\infty))$ is a cut off function defined as, $0\leq\rho \leq1$ such that $\rho=1$ in $[0,1]$ and $\rho=0 $ in $[2,\infty)$\B{this sentence with all the conditions is very long and hard to follow, especially which conditions are associated with which function. Maybe better split it into several sentences.}.
In the weak formulation of the microscopic problem, we choose $\phi=\es\phi_1$ as a test and obtain the identity
 \begin{multline}\label{eq:phi1weakform}
   \underbrace{ \es \rho_L\int_0^T\int_{\oll}  \dfrac{\partial^2u_L^\es}{\partial t^2} (\alpha_1(t)\alpha_2(\frac{t}{\es})\beta(\bar x) \psi_L(\frac{\bar x}{\es})\rho(\frac{x_1}{\es}))\di{x}\di{t}}_{I_1}\\
    +  \underbrace{\int_0^T\int_{\omm}  \dfrac{\partial}{\partial t}\left({\rho_M^\es} \dfrac{\partial u_M^\es}{\partial t}\right) ( \alpha_1(t)\alpha_2(\frac{t}{\es})\beta(\bar x) \psi_M(\frac{x}{\es}))\di{x}\di{t}}_{I_2} \\
    +   \underbrace{ \es \rho_R\int_0^T\int_{\orr} \dfrac{\partial^2u_R^\es}{\partial t^2} ( \alpha_1(t)\alpha_2(\frac{t}{\es})\beta(\bar x) \psi_R(\frac{\bar x}{\es})\rho(\frac{-x_1}{\es}))\di{x}\di{t}}_{I_3}\\
    +\underbrace{\es D_L\int_0^T\int_{\oll}\nabla u_L^\es \cdot \nabla ( \alpha_1(t)\alpha_2(\frac{t}{\es})\beta(\bar x) \psi_L(\frac{\bar x}{\es})\rho(\frac{x_1}{\es}))   \di{x}\di{t}}_{I_4}\\
    + \underbrace{\es D_R\int_0^T\int_{\orr} \nabla u_R^\es\cdot\nabla ( \alpha_1(t)\alpha_2(\frac{t}{\es})\beta(\bar x) \psi_R(\frac{\bar x}{\es})\rho(\frac{-x_1}{\es})) \di{x}\di{t}}_{I_5} \\
    +\underbrace{ \int_0^T\int_{\omm}  D_M^\es  \nabla u_M^\es\cdot\nabla (\alpha_1(t)\alpha_2(\frac{t}{\es})\beta(\bar x) \psi_M(\frac{x}{\es})) \di{x}\di{t}}_{I_6}\\
    = \underbrace{\es \int_0^T\int_{\oll} f_L(\alpha_1(t)\alpha_2(\frac{t}{\es})\beta(\bar x) \psi_L(\frac{\bar x}{\es})\rho(\frac{x_1}{\es})) \di{x}\di{t}}_{I_7}.
\end{multline}
We study the limit of each term in the aforementioned identity separately.
Using integration by parts with respect to the time variable, we can write the term $ I_1$ as
\begin{multline*} 
    I_1=-\es \rho_L\int_0^T\int_{\oll}  \dfrac{\partial u_L^\es}{\partial t} (\alpha_1'(t)\alpha_2(\frac{t}{\es})\beta(\bar x) \psi_L(\frac{\bar x}{\es})\rho(\frac{x_1}{\es}))\di{x}\di{t}\\
    -\rho_L\int_0^T\int_{\oll}  \dfrac{\partial u_L^\es}{\partial t} (\alpha_1(t)\alpha_2'(\frac{t}{\es})\beta(\bar x) \psi_L(\frac{\bar x}{\es})\rho(\frac{x_1}{\es}))\di{x}\di{t}.
\end{multline*}
The terms in the first integral are uniformly bounded with respect to the $L^2$-norm; therefore, using the Cauchy-Schwarz inequality, the first term converges to $0$ as $\es\rightarrow0$. 
%
For the second term of $I_1$, we have 

\begin{align*}
    \rho_L\int_0^T\int_{\oll}  \dfrac{\partial u_L^\es}{\partial t} (\alpha_1(t)\alpha_2'(\frac{t}{\es})\beta(\bar x) \psi_L(\frac{\bar x}{\es})\rho(\frac{x_1}{\es}))\di{x}\di{t}&\leq C \int_0^T\int_{\oll} \left| \dfrac{\partial u_L^\es}{\partial t}\right| |\rho(\frac{x_1}{\es}))| \di{x}\di{t}\\
    &\leq C  \left\|\ \dfrac{\partial u_L^\es}{\partial t}\right\|_{L^{2}(0,T;L^2(\oll))}T\Big(\int_0^{2\es}\rho(\frac{x_1}{\es})^2\di{x_1}\Big)^{\frac{1}{2}}\\
    &\leq \sqrt{\es} C.
\end{align*}
Thus, the second term also converges to $0$ as $\es\rightarrow 0$, this leads to $I_1\rightarrow0 $ as $\es\rightarrow 0$. Similarly, $I_3\rightarrow0 $ as $\es\rightarrow 0$.

Next, we look at the term $I_4$, and observe that
\begin{multline*}
  I_4
    =\es D_L\int_0^T\int_{\oll}\nabla u_L^\es \cdot \alpha_1(t)\alpha_2(\frac{t}{\es})  \nabla_{ x} \beta(\bar x) \psi_L(\frac{\bar x}{\es})\rho(\frac{x_1}{\es})   \di{x}\di{t}\\
    +D_L\int_0^T\int_{\oll}\nabla u_L^\es \cdot \alpha_1(t)\alpha_2(\frac{t}{\es})  \beta(\bar x)  \nabla_{ x}\psi_L(\frac{\bar x}{\es})\rho(\frac{x_1}{\es})   \di{x}\di{t}\\
    +D_L\int_0^T\int_{\oll}\nabla u_L^\es \cdot \alpha_1(t)\alpha_2(\frac{t}{\es})  \beta(\bar x)  \psi_L(\frac{\bar x}{\es})\rho'(\frac{x_1}{\es})e_1   \di{x}\di{t}.
\end{multline*}
it follows that the first term converges to $0$ as $\es\rightarrow0$. For the second term, we have
\begin{align*}
    \left|D_L\int_0^T\int_{\oll}\nabla u_L^\es \cdot \alpha_1(t)\alpha_2(\frac{t}{\es})  \beta(\bar x)  \nabla_{ x}\psi_L(\frac{\bar x}{\es})\rho(\frac{x_1}{\es})   \di{x}\di{t}\right|&\leq C \|\nabla u_L^\es\|_{L^2(0,T;L^2(\oll))}T\Big(\int_0^{2\es}\rho(\frac{x_1}{\es})^2\di{x_1})\Big)^{\frac{1}{2}}\\
    &\leq \sqrt{\es} C,
\end{align*}
where $C$ is independent of $\es$, and for the third term, we have
\begin{align*}
    \left|D_L\int_0^T\int_{\oll}\nabla u_L^\es \cdot \alpha_1(t)\alpha_2(\frac{t}{\es})  \beta(\bar x)  \psi_L(\frac{\bar x}{\es})\rho'(\frac{x_1}{\es}) e_1  \di{x}\di{t}\right|&\leq C \|\nabla u_L^\es\|_{L^2(0,T;L^2(\oll))}T\Big(\int_0^{2\es}\rho'(\frac{x_1}{\es})^2\di{x_1})\Big)^{\frac{1}{2}}\\
    &\leq \sqrt{\es} C,
\end{align*}
where $C$ is independent of $\es$. Consequently,  as $\es\rightarrow 0,$  we get $I_4\rightarrow 0.$ In a similar way, we deduce $I_5\rightarrow 0$ as $\es\rightarrow 0$. For the term $I_7$, each term inside the integral is bounded; therefore, as $\es\rightarrow 0$, $I_7\rightarrow 0$.

Now, we look at the term $I_2$ and we have 
\begin{multline*}
    I_2=- \int_0^T\int_{\omm}  {\rho_M^\es} \dfrac{\partial u_M^\es}{\partial t} ( \alpha_1'(t)\alpha_2(\frac{t}{\es})\beta(\bar x) \psi_M(\frac{x}{\es}))\di{x}\di{t}\\
    -\frac{1}{\es} \int_0^T\int_{\omm} {\rho_M^\es} \dfrac{\partial u_M^\es}{\partial t} ( \alpha_1(t)\alpha_2'(\frac{t}{\es})\beta(\bar x) \psi_M(\frac{x}{\es}))\di{x}\di{t}.
\end{multline*}
Relying on  \Cref{def:multiscale_convergence} and \eqref{eq:multiscale_convergence_for_time_derivative_term}, we arrive at
\begin{align*}
    \lim_{\es\rightarrow 0}\int_0^T\int_{\omm}  {\rho_M^\es} \dfrac{\partial u_M^\es}{\partial t} ( \alpha_1'(t)&\alpha_2(\frac{t}{\es})\beta(\bar x) \psi_M(\frac{x}{\es}))\di{x}\di{t}\\
    &=\lim_{\es\rightarrow 0}\es\cdot \lim_{\es\rightarrow 0} \frac 1\es \int_0^T\int_{\omm}  {\rho_M^\es} \dfrac{\partial u_M^\es}{\partial t} ( \alpha_1'(t)\alpha_2(\frac{t}{\es})\beta(\bar x) \psi_M(\frac{x}{\es}))\di{x}\di{t}\\
    &= 0,
\end{align*}
and
\begin{multline*}
    \lim_{\es\rightarrow 0}\frac{1}{\es} \int_0^T\int_{\omm} {\rho_M^\es} \dfrac{\partial u_M^\es}{\partial t} ( \alpha_1(t)\alpha_2'(\frac{t}{\es})\beta(\bar x) \psi_M(\frac{x}{\es}))\di{x}\di{t}\\
    = \mint  {\rho_M(t,\overline{x},\tau,y,z)} \left(\dfrac{\partial u_M^0(t,x)}{\partial t}+ \dfrac{\partial u_M^1(t,\overline{x},\tau,y)}{\partial \tau}\right) \\ ( \alpha_1(t)\alpha_2'(\tau)\beta(\bar x) \psi_M(y))\dzdy.
\end{multline*}
% By integration by parts with respect to the $\tau $ variable, we get
% \begin{multline*}
%      \lim_{\es\rightarrow 0}\frac{1}{\es} \int_0^T\int_{\omm} {\rho_M^\es} \dfrac{\partial u_M^\es}{\partial t} ( \alpha_1(t)\alpha_2'(\frac{t}{\es})\beta(\bar x) \psi_M(\frac{x}{\es}))\di{x}\di{t}\\
%     = -\int_0^T\int_{\Sigma}\int_\T\int_{Y}\int_Z \dfrac{\p}{\p \tau}\left( {\rho_M(t,\overline{x},\tau,y,z)} \left(\dfrac{\partial u_M^0(t,x)}{\partial t}+ \dfrac{\partial u_M^1(t,\overline{x},\tau,y)}{\partial \tau}\right)\right) \\ ( \alpha_1(t)\alpha_2(\tau)\beta(\bar x) \psi_M(y))\di{z}\di{y}\di{\tau}\di{x}\di{t}.
% \end{multline*}
It remains to show the limit of the term $I_6$. It holds that
\begin{multline*}
    I_6=\int_0^T\int_{\omm}  D_M^\es  \nabla u_M^\es\cdot (\alpha_1(t)\alpha_2(\frac{t}{\es})\nabla_{ x}\beta(\bar x) \psi_M(\frac{x}{\es})) \di{x}\di{t}\\
    +\frac1\es\int_0^T\int_{\omm}  D_M^\es  \nabla u_M^\es\cdot (\alpha_1(t)\alpha_2(\frac{t}{\es})\beta(\bar x) \nabla\psi_M(\frac{x}{\es})) \di{x}\di{t}.
\end{multline*}
Using \Cref{def:multiscale_convergence} and \eqref{eq:multiscale_convergence_for_gradient_term}, for the first term we have 
\begin{align*}
    \lim_{\es\rightarrow 0} \int_0^T\int_{\omm}  D_M^\es  \nabla u_M^\es\cdot (\alpha_1(t)&\alpha_2(\frac{t}{\es})\nabla\beta(\bar x)\psi_M(\frac{x}{\es})) \di{x}\di{t}\\
    &=\lim_{\es\rightarrow 0}\es\cdot \lim_{\es\rightarrow 0} \frac1\es \int_0^T\int_{\omm}  D_M^\es  \nabla u_M^\es\cdot (\alpha_1(t)\alpha_2(\frac{t}{\es})\nabla\beta(\bar x)\psi_M(\frac{x}{\es})) \di{x}\di{t}\\
    &= 0,
\end{align*}
and for the second term, we have 
\begin{multline*}
    \lim_{\es\rightarrow 0}\frac1\es\int_0^T\int_{\omm}  D_M^\es  \nabla u_M^\es\cdot (\alpha_1(t)\alpha_2(\frac{t}{\es})\beta(\bar x) \nabla\psi_M(\frac{x}{\es})) \di{x}\di{t}\\
    =\int_0^T\int_{\Sigma}\int_\T\int_{Y}\int_Z   D_M(t,\overline{x},\tau,y,z)
    \left( \nabla_{\overline{x}}u_M^0(t,\overline{x})+\nabla_y u_M^1(t,\overline{x},\tau,y)+\nabla_zu_M^2(t,\overline{x},\tau,y,z)\right)\\
    \cdot ( \alpha_1(t)\alpha_2(\tau)\beta(\bar x) \nabla_y\psi_M(y)) \di{z}\di{y}\di{\tau}\di{x}\di{t}.
\end{multline*}

% Now doing integration by parts with respect to the $y$ variable, we get
% \begin{multline*}
%     \lim_{\es\rightarrow 0}\frac1\es\int_0^T\int_{\omm}  D_M^\es  \nabla u_M^\es\cdot (\alpha_1(t)\alpha_2(\frac{t}{\es})\beta(\bar x) \nabla\psi_M(\frac{x}{\es})) \di{x}\di{t}\\
%     =-\int_0^T\int_{\Sigma}\int_\T\int_{Y}\int_Z  \nabla_y\cdot\bigg( D_M(t,\overline{x},\tau,y,z)\\
%     \left( \nabla_{\overline{x}}u_M^0(t,\overline{x})+\nabla_y u_M^1(t,\overline{x},\tau,y)+\nabla_zu_M^2(t,\overline{x},\tau,y,z)\right)\bigg)\\
%     \cdot ( \alpha_1(t)\alpha_2(\tau)\beta(\bar x) \psi_M(y)) \di{z}\di{y}\di{\tau}\di{x}\di{t}.
% \end{multline*}
 

Collecting the limits of $I_1-I_7$, as $\es\rightarrow 0$, from the identity \eqref{eq:phi1weakform} we arrive at
\begin{multline*}
     -\int_0^T\int_{\Sigma}\int_\T\int_{Y}\int_Z  {\rho_M(t,\overline{x},\tau,y,z)} \left(\dfrac{\partial u_M^0(t,x)}{\partial t}+ \dfrac{\partial u_M^1(t,\overline{x},\tau,y)}{\partial \tau}\right) \\ ( \alpha_1(t)\alpha_2'(\tau)\beta(\bar x) \psi_M(y))\di{z}\di{y}\di{\tau}\di{x}\di{t}\\
     +\int_0^T\int_{\Sigma}\int_\T\int_{Y}\int_Z   D_M(t,\overline{x},\tau,y,z)
    \left( \nabla_{\overline{x}}u_M^0(t,\overline{x})+\nabla_y u_M^1(t,\overline{x},\tau,y)+\nabla_zu_M^2(t,\overline{x},\tau,y,z)\right)\\
    \cdot ( \alpha_1(t)\alpha_2(\tau)\beta(\bar x) \nabla_y\psi_M(y)) \di{z}\di{y}\di{\tau}\di{x}\di{t}=0.
\end{multline*}
% Now, to obtain a strong form, we perform the integration by parts with respect to the $\tau$ variable on the first term and integration by parts with respect to  the $y$ variable in the second term, and we obtain the relation
% \begin{multline*}
%     \int_0^T\int_{\Sigma}\int_\T\int_{Y}\int_Z \dfrac{\p}{\p \tau}\left( {\rho_M(t,\overline{x},\tau,y,z)} \left(\dfrac{\partial u_M^0(t,x)}{\partial t}+ \dfrac{\partial u_M^1(t,\overline{x},\tau,y)}{\partial \tau}\right)\right) \\ ( \alpha_1(t)\alpha_2(\tau)\beta(\bar x) \psi_M(y))\di{z}\di{y}\di{\tau}\di{x}\di{t}\\
%     -\int_0^T\int_{\Sigma}\int_\T\int_{Y}\int_Z  \nabla_y\cdot\bigg( D_M(t,\overline{x},\tau,y,z)\\
%     \left( \nabla_{\overline{x}}u_M^0(t,\overline{x})+\nabla_y u_M^1(t,\overline{x},\tau,y)+\nabla_zu_M^2(t,\overline{x},\tau,y,z)\right)\bigg)\\
%     \cdot ( \alpha_1(t)\alpha_2(\tau)\beta(\bar x) \psi_M(y)) \di{z}\di{y}\di{\tau}\di{x}\di{t}=0
% \end{multline*}

That is the weak form of the identity \eqref{eq:strong_form_u_1} and \eqref{eq:strong_form_u_1_bc}. 
% \begin{multline}
%    \dfrac{\p}{\p \tau} \int_Z \left( {\rho_M(t,\overline{x},\tau,y,z)} \left(\dfrac{\partial u_M^0(t,x)}{\partial t}+ \dfrac{\partial u_M^1(t,\overline{x},\tau,y)}{\partial \tau}\right)\right)\di{z}\\
%      -\mathrm{div_y}\int_Z \bigg( D_M(t,\overline{x},\tau,y,z)
%     \left( \nabla_{\overline{x}}u_M^0(t,\overline{x})+\nabla_y u_M^1(t,\overline{x},\tau,y)+\nabla_zu_M^2(t,\overline{x},\tau,y,z)\right)\bigg)\di{z}=0\\
%     \mbox{for a.e.} (t,\overline{x},\tau,y)\in (0,T)\times\Sigma\times (0,1)\times Y.
% \end{multline}
% We can write the strong form of the aforementioned identity as 
% \begin{align}
%     \partial_{\tau}\left(\rho(t,\overline{x},\tau, y,z)\partial_{\tau}u_M^1\right)+\partial_{\tau}\rho(t,\overline{x},\tau, y,z)\partial_t u_M^0\hspace{1.5cm}\qquad&\\
%  -\mathrm{div}_y\left(D(t,\overline{x},\tau, y,z)\left(\nabla_{\overline{x}}u_M^0+\nabla_y u_M^1 +\nabla_zu_M^2\right)\right)&=0
% \end{align}

We are left to show the identity \eqref{eq:strong_form_for_u_2}. To do so, we construct a new test function $\phi_2$ in $(0,T)\times\Omega$.
Let
\begin{align*}
    \phi_2:=\begin{cases} 
      \alpha_1(t)\alpha_2(\frac{t}{\es})\beta(\bar x) \psi_L(\frac{\bar x}{\es})\rho(\frac{x_1}{\es})\gamma(\frac{x_2}{\es^2}) &  \qquad \mbox{in}\qquad (0,T)\times \oll\\
      \alpha_1(t)\alpha_2(\frac{t}{\es})\beta(\bar x) \psi_M(\frac{x}{\es})\gamma(\frac{x_2}{\es^2}) & \qquad \mbox{in}\qquad (0,T)\times \omm\\
      \alpha_1(t)\alpha_2(\frac{t}{\es})\beta(\bar x) \psi_R(\frac{\bar x}{\es})\rho(\frac{-x_1}{\es})\gamma(\frac{x_2}{\es^2}) & \qquad \mbox{in}\qquad (0,T)\times \orr,
   \end{cases}
\end{align*}
where $\alpha_1,\alpha_2,\beta, \psi_L,\psi_M,\psi_R,\rho$ are functions defined as before and $\gamma\in C_c^{\infty}([0,1])$ which extends periodically to whole $\mathbb{R}$. Choosing this well-constructed test function $\phi=\es^2 \phi_2$ in the weak formulation \eqref{eq:weakform}, we get
 \begin{multline}\label{eq:phi2weakform}
   \underbrace{ \es^2 \rho_L\int_0^T\int_{\oll}  \dfrac{\partial^2u_L^\es}{\partial t^2} (\alpha_1(t)\alpha_2(\frac{t}{\es})\beta(\bar x) \psi_L(\frac{\bar x}{\es})\rho(\frac{x_1}{\es})\gamma(\frac{x_2}{\es^2}))\di{x}\di{t}}_{I_8}\\
    +  \underbrace{\es\int_0^T\int_{\omm}  \dfrac{\partial}{\partial t}\left({\rho_M^\es} \dfrac{\partial u_M^\es}{\partial t}\right) ( \alpha_1(t)\alpha_2(\frac{t}{\es})\beta(\bar x) \psi_M(\frac{x}{\es})\gamma(\frac{x_2}{\es^2}))\di{x}\di{t}}_{I_9} \\
    +   \underbrace{ \es^2 \rho_R\int_0^T\int_{\orr} \dfrac{\partial^2u_R^\es}{\partial t^2} ( \alpha_1(t)\alpha_2(\frac{t}{\es})\beta(\bar x) \psi_R(\frac{\bar x}{\es})\rho(\frac{-x_1}{\es})\gamma(\frac{x_2}{\es^2}))\di{x}\di{t}}_{I_{10}}\\
    +\underbrace{\es^2 D_L\int_0^T\int_{\oll}\nabla u_L^\es \cdot \nabla ( \alpha_1(t)\alpha_2(\frac{t}{\es})\beta(\bar x) \psi_L(\frac{\bar x}{\es})\rho(\frac{x_1}{\es})\gamma(\frac{x_2}{\es^2}))   \di{x}\di{t}}_{I_{11}}\\
    + \underbrace{\es^2 D_R\int_0^T\int_{\orr} \nabla u_R^\es\cdot\nabla ( \alpha_1(t)\alpha_2(\frac{t}{\es})\beta(\bar x) \psi_R(\frac{\bar x}{\es})\rho(\frac{-x_1}{\es})\gamma(\frac{x_2}{\es^2})) \di{x}\di{t}}_{I_{12}} \\
    +\underbrace{ \es\int_0^T\int_{\omm}  D_M^\es  \nabla u_M^\es\cdot\nabla (\alpha_1(t)\alpha_2(\frac{t}{\es})\beta(\bar x) \psi_M(\frac{x}{\es})\gamma(\frac{x_2}{\es^2})) \di{x}\di{t}}_{I_{13}}\\
    = \underbrace{\es^2 \int_0^T\int_{\oll} f_L(\alpha_1(t)\alpha_2(\frac{t}{\es})\beta(\bar x) \psi_L(\frac{\bar x}{\es})\rho(\frac{x_1}{\es})\gamma(\frac{x_2}{\es^2})) \di{x}\di{t}}_{I_{14}}.
\end{multline}
Similar to the case of \eqref{eq:phi1weakform}, we derive that $I_8,I_9,I_{10},I_{11},I_{12},I_{14}\rightarrow 0$ as $\es\rightarrow 0$. It remains to check the limit of the term $I_{13}$ as $\es\rightarrow 0$.  It holds that
\begin{multline*}
    I_{13}= \es\int_0^T\int_{\omm}  D_M^\es  \nabla u_M^\es\cdot (\alpha_1(t)\alpha_2(\frac{t}{\es})\nabla\beta(\bar x) \psi_M(\frac{x}{\es})\gamma(\frac{x_2}{\es^2})) \di{x}\di{t}\\
    +\int_0^T\int_{\omm}  D_M^\es  \nabla u_M^\es\cdot (\alpha_1(t)\alpha_2(\frac{t}{\es})\beta(\bar x) \nabla\psi_M(\frac{x}{\es})\gamma(\frac{x_2}{\es^2})) \di{x}\di{t}\\
    +\frac1\es\int_0^T\int_{\omm}  D_M^\es  \nabla u_M^\es\cdot (\alpha_1(t)\alpha_2(\frac{t}{\es})\beta(\bar x) \psi_M(\frac{x}{\es})\gamma'(\frac{x_2}{\es^2})e_2) \di{x}\di{t}.
\end{multline*}
By means of \Cref{def:multiscale_convergence} and \eqref{eq:multiscale_convergence_for_gradient_term}, we see that the first two terms converge to $0$. For the final term, taking the limit $\es\rightarrow 0$, we obtain
\begin{align*}
    &\!\!\!\!\lim_{\es\rightarrow 0}\frac1\es\int_0^T\int_{\omm}  D_M^\es  \nabla u_M^\es\cdot (\alpha_1(t)\alpha_2(\frac{t}{\es})\beta(\bar x) \psi_M(\frac{x}{\es})\gamma'(\frac{x_2}{\es^2})) e_2\di{x}\di{t}\\
    &=\int_0^T\int_{\Sigma}\int_\T\int_{Y}\int_Z   D_M(t,\overline{x},\tau,y,z)
    \left( \nabla_{\overline{x}}u_M^0(t,\overline{x})+\nabla_y u_M^1(t,\overline{x},\tau,y)+\nabla_zu_M^2(t,\overline{x},\tau,y,z)\right)\\
    &\qquad\qquad\qquad\qquad \cdot ( \alpha_1(t)\alpha_2(\tau)\beta(\bar x) \psi_M(y)\gamma'(z)e_2) \di{z}\di{y}\di{\tau}\di{x}\di{t}.
\end{align*}
Therefore, taking the limit $\es\rightarrow 0$ in \eqref{eq:phi2weakform}, we deduce the following identity
\begin{multline*}
    \int_0^T\int_{\Sigma}\int_\T\int_{Y}\int_Z   D_M(t,\overline{x},\tau,y,z)
    \left( \nabla_{\overline{x}}u_M^0(t,\overline{x})+\nabla_y u_M^1(t,\overline{x},\tau,y)+\nabla_zu_M^2(t,\overline{x},\tau,y,z)\right)\\
    \cdot ( \alpha_1(t)\alpha_2(\tau)\beta(\bar x) \psi_M(y)\gamma'(z)e_2) \di{z}\di{y}\di{\tau}\di{x}\di{t}=0.
\end{multline*}
% Now, to obtain a strong form, we do integration by parts with respect to the $z$ variable and  
% obtain
% \begin{multline*}
%     \int_0^T\int_{\Sigma}\int_\T\int_{Y}\int_Z \nabla_z\cdot \Big(   D_M(t,\overline{x},\tau,y,z)
%     \left( \nabla_{\overline{x}}u_M^0(t,\overline{x})+\nabla_y u_M^1(t,\overline{x},\tau,y)+\nabla_zu_M^2(t,\overline{x},\tau,y,z)\right)\Big)\\
%     \cdot ( \alpha_1(t)\alpha_2(\tau)\beta(\bar x) \psi_M(y)\gamma(z)) \di{z}\di{y}\di{\tau}\di{x}\di{t}=0.
% \end{multline*}
This is precisely the weak form of  \eqref{eq:strong_form_for_u_2}.
% \begin{multline}
%     -\mathrm{div_z}\Big(   D_M(t,\overline{x},\tau,y,z)
%     \left( \nabla_{\overline{x}}u_M^0(t,\overline{x})+\nabla_y u_M^1(t,\overline{x},\tau,y)+\nabla_zu_M^2(t,\overline{x},\tau,y,z)\right)\Big)=0\\
%     \mbox{for a.e. } (t,\overline{x},\tau,y,z)\in (0,T)\times\Sigma\times (0,1)\times Y\times Z.
% \end{multline}
\end{proof}

\end{document}
